% answers.tex —— 2004 年·普通高等学校招生考试·江苏卷 · 参考答案与详细解析片段
% 由 tools/generate_projects.py 依来源总账自 2004/2004_数学真题_江苏_试题.md + 2004/2004_数学真题_江苏_解析正文.md + 2004/2004_数学真题_江苏_答案要点.md 生成，勿手改（改源 MD 后重跑）


\begin{gktip}
固定源 PDF 为印刷页码 178–179 的两页题面，不附答案。以下解析均依据固定源独立推导复算；第 6 题条形图与第 18 题立体几何图在 LaTeX 工程中以 TikZ 矢量源重绘。第 9 题题面漏印文字经公开权威资料考证补全。
\end{gktip}

\begin{answer}
\gksolhead{1}{}

\textbf{题目}：

设集合 $P = \{1, 2, 3, 4\}$，$Q = \{x \mid |x| \leqslant 2, x \in \mathbf{R} \}$，则 $P \cap Q$ 等于\mbox{（　　）}

\begin{gkopts}(2)
\task A．$\{1, 2\}$
\task B．$\{3, 4\}$
\task C．$\{1\}$
\task D．$\{-2, -1, 0, 1, 2\}$
\end{gkopts}

\Needspace{4\baselineskip}
\textbf{解析}

由集合 $Q = \{x \mid |x| \leqslant 2, x \in \mathbf{R} \}$ 得 $-2 \leqslant x \leqslant 2$．
集合 $P = \{1, 2, 3, 4\}$，
所以 $P \cap Q = \{1, 2\}$．
故选 A．

\gkanswer{答案：A}

\end{answer}
\begin{answer}
\gksolhead{2}{}

\textbf{题目}：

函数 $y = 2\cos^2 x + 1$（$x \in \mathbf{R}$）的最小正周期为\mbox{（　　）}

\begin{gkopts}(4)
\task A．$\dfrac{\pi}{2}$
\task B．$\pi$
\task C．$2\pi$
\task D．$4\pi$
\end{gkopts}

\Needspace{4\baselineskip}
\textbf{解析}

\Needspace{5\baselineskip}
利用二倍角余弦公式降幂：
\nopagebreak[4]
\[
y = 2\cos^2 x + 1 = (\cos 2x + 1) + 1 = \cos 2x + 2.
\]
\Needspace{5\baselineskip}
其最小正周期为
\nopagebreak[4]
\[
T = \frac{2\pi}{2} = \pi.
\]
故选 B．

\gkanswer{答案：B}

\end{answer}
\begin{answer}
\gksolhead{3}{}

\textbf{题目}：

从 4 名男生和 3 名女生中选出 4 人参加某个座谈会，若这 4 人中必须既有男生又有女生，则不同的选法共有\mbox{（　　）}

\begin{gkopts}(4)
\task A．$140$ 种
\task B．$120$ 种
\task C．$35$ 种
\task D．$34$ 种
\end{gkopts}

\Needspace{4\baselineskip}
\textbf{解析}

\Needspace{5\baselineskip}
从 4 名男生和 3 名女生共 7 人中任选 4 人，不同的选法总数为
\nopagebreak[4]
\[
\binom{7}{4} = \frac{7 \times 6 \times 5}{3 \times 2 \times 1} = 35.
\]
若选出的 4 人全为男生，选法有 $\binom{4}{4} = 1$ 种；
因为女生只有 3 人，所以不可能 4 人全为女生．
\Needspace{5\baselineskip}
因此，4 人中既有男生又有女生的选法共有
\nopagebreak[4]
\[
35 - 1 = 34 \text{ \text{种}}.
\]
故选 D．

\gkanswer{答案：D}

\end{answer}
\begin{answer}
\gksolhead{4}{}

\textbf{题目}：

一平面截一球得到直径是 $6\text{ cm}$ 的圆面，球心到这个平面的距离是 $4\text{ cm}$，则该球的体积是\mbox{（　　）}

\begin{gkopts}(2)
\task A．$\dfrac{100\pi}{3}\text{ cm}^3$
\task B．$\dfrac{208\pi}{3}\text{ cm}^3$
\task C．$\dfrac{500\pi}{3}\text{ cm}^3$
\task D．$\dfrac{416\sqrt{3}\pi}{3}\text{ cm}^3$
\end{gkopts}

\Needspace{4\baselineskip}
\textbf{解析}

设截面圆半径为 $r$，球半径为 $R$，球心到截面的距离为 $d$．
由题意，截面圆直径为 $6\text{ cm}$，故 $r = 3\text{ cm}$，且 $d = 4\text{ cm}$．
\Needspace{5\baselineskip}
由勾股定理可得球半径
\nopagebreak[4]
\[
R = \sqrt{r^2 + d^2} = \sqrt{3^2 + 4^2} = 5\text{ cm}.
\]
\Needspace{5\baselineskip}
因此该球的体积为
\nopagebreak[4]
\[
V = \frac{4}{3}\pi R^3 = \frac{4}{3}\pi \times 5^3 = \frac{500\pi}{3}\text{ cm}^3.
\]
故选 C．

\gkanswer{答案：C}

\end{answer}
\begin{answer}
\gksolhead{5}{}

\textbf{题目}：

若双曲线 $\dfrac{x^2}{8} - \dfrac{y^2}{b^2} = 1$ 的一条准线与抛物线 $y^2 = 8x$ 的准线重合，则双曲线离心率为\mbox{（　　）}

\begin{gkopts}(4)
\task A．$\sqrt{2}$
\task B．$2\sqrt{2}$
\task C．$4$
\task D．$4\sqrt{2}$
\end{gkopts}

\Needspace{4\baselineskip}
\textbf{解析}

抛物线 $y^2 = 8x = 2px$ 的准线方程为 $x = -\dfrac{p}{2} = -2$．
双曲线 $\dfrac{x^2}{8} - \dfrac{y^2}{b^2} = 1$ 的焦点在 $x$ 轴上，其中 $a^2 = 8$，其左准线方程为 $x = -\dfrac{a^2}{c}$．
\Needspace{5\baselineskip}
由题意准线重合，故
\nopagebreak[4]
\[
-\frac{8}{c} = -2 \implies c = 4.
\]
\Needspace{5\baselineskip}
因此双曲线的离心率为
\nopagebreak[4]
\[
e = \frac{c}{a} = \frac{4}{\sqrt{8}} = \frac{4}{2\sqrt{2}} = \sqrt{2}.
\]
故选 A．

\gkanswer{答案：A}

\end{answer}
\begin{answer}
\gksolhead{6}{}

\textbf{题目}：

某校为了了解学生的课外阅读情况，随机调查了 50 名学生，得到他们在某一天各自课外阅读所用时间的数据，结果用下面的条形图表示．根据条形图可得这 50 名学生这一天平均每人的课外阅读时间为\mbox{（　　）}
\nopagebreak[4]
\begin{center}\begin{tikzpicture}[>=stealth, xscale=1.5, yscale=0.18]
  % 网格虚线
  \draw[dashed, gray!50] (0, 5) -- (4.5, 5);
  \draw[dashed, gray!50] (0, 10) -- (4.5, 10);
  \draw[dashed, gray!50] (0, 15) -- (4.5, 15);
  \draw[dashed, gray!50] (0, 20) -- (4.5, 20);

  % 条形柱
  \draw[fill=gray!30] (-0.15, 0) rectangle (0.15, 5);
  \draw[fill=gray!30] (0.85, 0) rectangle (1.15, 20);
  \draw[fill=gray!30] (1.85, 0) rectangle (2.15, 10);
  \draw[fill=gray!30] (2.85, 0) rectangle (3.15, 10);
  \draw[fill=gray!30] (3.85, 0) rectangle (4.15, 5);

  % 坐标轴
  \draw[->, thick] (-0.3, 0) -- (4.8, 0) node[right] {\small 时间\,(小时)};
  \draw[->, thick] (0, 0) -- (0, 23) node[above] {\small 人数\,(人)};

  % 纵轴刻度
  \node[left] at (0, 0) {\small $0$};
  \node[left] at (0, 5) {\small $5$};
  \node[left] at (0, 10) {\small $10$};
  \node[left] at (0, 15) {\small $15$};
  \node[left] at (0, 20) {\small $20$};

  % 横轴刻度与标签
  \node[below] at (0, 0) {\small $0$};
  \node[below] at (1, 0) {\small $0.5$};
  \node[below] at (2, 0) {\small $1.0$};
  \node[below] at (3, 0) {\small $1.5$};
  \node[below] at (4, 0) {\small $2.0$};
\end{tikzpicture}
\end{center}
\nopagebreak[4]
\begin{gkopts}(4)
\task A．$0.6$ 小时
\task B．$0.9$ 小时
\task C．$1.0$ 小时
\task D．$1.5$ 小时
\end{gkopts}

\Needspace{4\baselineskip}
\textbf{解析}
\nopagebreak[4]
\begin{center}\begin{tikzpicture}[>=stealth, xscale=1.5, yscale=0.18]
  % 网格虚线
  \draw[dashed, gray!50] (0, 5) -- (4.5, 5);
  \draw[dashed, gray!50] (0, 10) -- (4.5, 10);
  \draw[dashed, gray!50] (0, 15) -- (4.5, 15);
  \draw[dashed, gray!50] (0, 20) -- (4.5, 20);

  % 条形柱
  \draw[fill=gray!30] (-0.15, 0) rectangle (0.15, 5);
  \draw[fill=gray!30] (0.85, 0) rectangle (1.15, 20);
  \draw[fill=gray!30] (1.85, 0) rectangle (2.15, 10);
  \draw[fill=gray!30] (2.85, 0) rectangle (3.15, 10);
  \draw[fill=gray!30] (3.85, 0) rectangle (4.15, 5);

  % 坐标轴
  \draw[->, thick] (-0.3, 0) -- (4.8, 0) node[right] {\small 时间\,(小时)};
  \draw[->, thick] (0, 0) -- (0, 23) node[above] {\small 人数\,(人)};

  % 纵轴刻度
  \node[left] at (0, 0) {\small $0$};
  \node[left] at (0, 5) {\small $5$};
  \node[left] at (0, 10) {\small $10$};
  \node[left] at (0, 15) {\small $15$};
  \node[left] at (0, 20) {\small $20$};

  % 横轴刻度与标签
  \node[below] at (0, 0) {\small $0$};
  \node[below] at (1, 0) {\small $0.5$};
  \node[below] at (2, 0) {\small $1.0$};
  \node[below] at (3, 0) {\small $1.5$};
  \node[below] at (4, 0) {\small $2.0$};
\end{tikzpicture}
\end{center}
\nopagebreak[4]
由条形图可知，50 名学生中各自课外阅读时间及对应人数如下：
\begin{itemize}
\item 阅读时间 0 小时：5 人；
\item 阅读时间 0.5 小时：20 人；
\item 阅读时间 1.0 小时：10 人；
\item 阅读时间 1.5 小时：10 人；
\item 阅读时间 2.0 小时：5 人．
\end{itemize}

\Needspace{5\baselineskip}
50 名学生课外阅读的总时间为
\nopagebreak[4]
\[
0 \times 5 + 0.5 \times 20 + 1.0 \times 10 + 1.5 \times 10 + 2.0 \times 5 = 0 + 10 + 10 + 15 + 10 = 45 \text{ \text{小时}}.
\]
\Needspace{5\baselineskip}
因此平均每人的课外阅读时间为
\nopagebreak[4]
\[
\frac{45}{50} = 0.9 \text{ \text{小时}}.
\]
故选 B．

\gkanswer{答案：B}

\end{answer}
\begin{answer}
\gksolhead{7}{}

\textbf{题目}：

$(2x + \sqrt{x})^4$ 的展开式中 $x^3$ 的系数是\mbox{（　　）}

\begin{gkopts}(4)
\task A．$6$
\task B．$12$
\task C．$24$
\task D．$48$
\end{gkopts}

\Needspace{4\baselineskip}
\textbf{解析}

\Needspace{5\baselineskip}
由二项式展开公式，$(2x + \sqrt{x})^4$ 的通项为
\nopagebreak[4]
\[
T_{r+1} = \binom{4}{r} (2x)^{4-r} (\sqrt{x})^r = \binom{4}{r} 2^{4-r} x^{4-r} \cdot x^{r/2} = \binom{4}{r} 2^{4-r} x^{4 - r/2}.
\]
令 $4 - \dfrac{r}{2} = 3$，解得 $r = 2$．
\Needspace{5\baselineskip}
此时 $x^3$ 的系数为
\nopagebreak[4]
\[
\binom{4}{2} 2^{4-2} = 6 \times 4 = 24.
\]
故选 C．

\gkanswer{答案：C}

\end{answer}
\begin{answer}
\gksolhead{8}{}

\textbf{题目}：

若函数 $y = \log_a(x + b)$（$a > 0, a \neq 1$）的图象过两点 $(-1, 0)$ 和 $(0, 1)$，则\mbox{（　　）}

\begin{gkopts}(2)
\task A．$a = 2, b = 2$
\task B．$a = \sqrt{2}, b = 2$
\task C．$a = 2, b = 1$
\task D．$a = \sqrt{2}, b = \sqrt{2}$
\end{gkopts}

\Needspace{4\baselineskip}
\textbf{解析}

\Needspace{5\baselineskip}
将点 $(-1, 0)$ 和 $(0, 1)$ 代入函数解析式 $y = \log_a(x + b)$：
\nopagebreak[4]
\[
\begin{cases}
\log_a(-1 + b) = 0, \\
\log_a(0 + b) = 1.
\end{cases}
\]
由第一式得 $-1 + b = 1 \implies b = 2$；
将 $b = 2$ 代入第二式得 $\log_a 2 = 1 \implies a = 2$．
因此 $a = 2, b = 2$．
故选 A．

\gkanswer{答案：A}

\end{answer}
\begin{answer}
\gksolhead{9}{}

\textbf{题目}：

将一颗质地均匀的骰子（它是一种各面上分别标有点数 1，2，3，4，5，6 的正方体玩具）先后抛掷 3 次，至少出现一次 6 点的概率是\mbox{（　　）}

\begin{gkopts}(2)
\task A．$\dfrac{5}{216}$
\task B．$\dfrac{25}{216}$
\task C．$\dfrac{31}{216}$
\task D．$\dfrac{91}{216}$
\end{gkopts}

\Needspace{4\baselineskip}
\textbf{解析}

\begin{gktip}
考证说明：原卷印刷文字漏印事件名称“一次 6 点”，完整题面为“至少出现一次 6 点的概率”．
\end{gktip}

掷一颗均匀骰子一次，出现 6 点的概率为 $\dfrac{1}{6}$，不出现 6 点的概率为 $\dfrac{5}{6}$．
\Needspace{5\baselineskip}
抛掷 3 次骰子，3 次均不出现 6 点的概率为
\nopagebreak[4]
\[
P(\text{\text{均不出现} 6 \text{点}}) = \left(\frac{5}{6}\right)^3 = \frac{125}{216}.
\]
\Needspace{5\baselineskip}
因此，“至少出现一次 6 点”的概率为对立事件的概率：
\nopagebreak[4]
\[
P = 1 - \frac{125}{216} = \frac{91}{216}.
\]
故选 D．

\gkanswer{答案：D}

\end{answer}
\begin{answer}
\gksolhead{10}{}

\textbf{题目}：

函数 $f(x) = x^3 - 3x + 1$ 在闭区间 $[-3, 0]$ 上的最大值、最小值分别是\mbox{（　　）}

\begin{gkopts}(4)
\task A．$1, -1$
\task B．$1, -17$
\task C．$3, -17$
\task D．$9, -19$
\end{gkopts}

\Needspace{4\baselineskip}
\textbf{解析}

\Needspace{5\baselineskip}
函数 $f(x) = x^3 - 3x + 1$，求导得
\nopagebreak[4]
\[
f'(x) = 3x^2 - 3 = 3(x - 1)(x + 1).
\]
在闭区间 $[-3, 0]$ 内，令 $f'(x) = 0$ 得驻点 $x = -1$．
\Needspace{5\baselineskip}
计算各关键点处的函数值：
\nopagebreak[4]
\[
f(-3) = (-3)^3 - 3(-3) + 1 = -27 + 9 + 1 = -17,
\]
\nopagebreak[4]
\[
f(-1) = (-1)^3 - 3(-1) + 1 = -1 + 3 + 1 = 3,
\]
\nopagebreak[4]
\[
f(0) = 0^3 - 3(0) + 1 = 1.
\]
比较可知，函数在 $[-3, 0]$ 上的最大值为 $3$，最小值为 $-17$．
故选 C．

\gkanswer{答案：C}

\end{answer}
\begin{answer}
\gksolhead{11}{}

\textbf{题目}：

设 $k > 1, f(x) = k(x - 1)$（$x \in \mathbf{R}$）．在平面直角坐标系 $xOy$ 中，函数 $y = f(x)$ 的图象与 $x$ 轴交于 $A$ 点，它的反函数 $y = f^{-1}(x)$ 的图象与 $y$ 轴交于 $B$ 点，并且这两个函数的图象交于 $P$ 点．已知四边形 $OAPB$ 的面积是 3，则 $k$ 等于\mbox{（　　）}

\begin{gkopts}(4)
\task A．$3$
\task B．$\dfrac{3}{2}$
\task C．$\dfrac{4}{3}$
\task D．$\dfrac{6}{5}$
\end{gkopts}

\Needspace{4\baselineskip}
\textbf{解析}

函数 $y = f(x) = k(x - 1)$（$k > 1$）：
令 $y = 0$，得 $x = 1$，故点 $A(1, 0)$．
由反函数的图象与原函数图象关于直线 $y = x$ 对称可知，反函数 $y = f^{-1}(x)$ 与 $y$ 轴的交点坐标为 $B(0, 1)$．
两函数的图象交于点 $P$，交点必在直线 $y = x$ 上：
联立 $x = k(x - 1) \implies x = kx - k \implies (k - 1)x = k$．
因为 $k > 1$，解得 $x = \dfrac{k}{k - 1}$，即 $P\left(\dfrac{k}{k-1}, \dfrac{k}{k-1}\right)$．
\Needspace{5\baselineskip}
四边形 $OAPB$ 的面积可沿对角线 $OP$ 分割为两个关于 $y=x$ 对称的三角形 $\triangle OAP$ 与 $\triangle OBP$：
\nopagebreak[4]
\[
S_{\triangle OAP} = \frac{1}{2} \times |OA| \times y_P = \frac{1}{2} \times 1 \times \frac{k}{k - 1} = \frac{k}{2(k - 1)}.
\]
\Needspace{5\baselineskip}
因此四边形 $OAPB$ 的面积为
\nopagebreak[4]
\[
S_{OAPB} = 2 S_{\triangle OAP} = \frac{k}{k - 1}.
\]
\Needspace{5\baselineskip}
已知面积为 3，故
\nopagebreak[4]
\[
\frac{k}{k - 1} = 3 \implies k = 3k - 3 \implies 2k = 3 \implies k = \frac{3}{2}.
\]
故选 B．

\gkanswer{答案：B}

\end{answer}
\begin{answer}
\gksolhead{12}{}

\textbf{题目}：

设函数 $f(x) = -\dfrac{x}{1 + |x|}$（$x \in \mathbf{R}$），区间 $M = [a, b]$（$a < b$），集合 $N = \{y \mid y = f(x), x \in M\}$，则使 $M = N$ 成立的实数对 $(a, b)$ 有\mbox{（　　）}

\begin{gkopts}(4)
\task A．$0$ 个
\task B．$1$ 个
\task C．$2$ 个
\task D．无数多个
\end{gkopts}

\Needspace{4\baselineskip}
\textbf{解析}

函数 $f(x) = -\dfrac{x}{1 + |x|}$（$x \in \mathbf{R}$）：
当 $x \geqslant 0$ 时，$f(x) = -\dfrac{x}{1 + x} = -1 + \dfrac{1}{1 + x}$，为严格单调递减函数；
当 $x < 0$ 时，$f(x) = -\dfrac{x}{1 - x} = 1 - \dfrac{1}{1 - x}$，亦为严格单调递减函数；
且 $f(-x) = -f(x)$ 为奇函数，故 $f(x)$ 在 $\mathbf{R}$ 上严格单调递减．
\Needspace{5\baselineskip}
当 $x \in M = [a, b]$（$a < b$）时，函数值域为
\nopagebreak[4]
\[
N = [f(b), f(a)].
\]
\Needspace{5\baselineskip}
要使 $M = N$，必须且只需
\nopagebreak[4]
\[
\begin{cases}
f(b) = a, \\
f(a) = b.
\end{cases}
\]
1. 若 $0 \leqslant a < b$，由单调性知 $f(b) < f(a) \leqslant 0$，值域包含于 $(-\infty, 0]$，而定义域包含于 $[0, +\infty)$，不可能满足 $M = N$；
2. 若 $a < b \leqslant 0$，同理值域包含于 $[0, +\infty)$，定义域包含于 $(-\infty, 0]$，不可能满足 $M = N$；
3. 若 $a < 0 < b$，则：
   由 $a < 0$ 得 $f(a) = -\dfrac{a}{1 - a}$，代入 $f(a) = b$ 得 $b = -\dfrac{a}{1 - a} \implies b - ab = -a \implies a + b = ab$；
   由 $b > 0$ 得 $f(b) = -\dfrac{b}{1 + b}$，代入 $f(b) = a$ 得 $a = -\dfrac{b}{1 + b} \implies a + ab = -b \implies a + b = -ab$．
   两式相减得 $2ab = 0$．但已知 $a < 0 < b$，必有 $ab < 0 \neq 0$，产生矛盾！
综上所述，不存在实数对 $(a, b)$ 使得 $M = N$ 成立，即实数对个数为 0 个．
故选 A．

\gkanswer{答案：A}

\end{answer}
\begin{answer}
\gksolhead{13}{}

\textbf{题目}：

二次函数 $y = ax^2 + bx + c$（$x \in \mathbf{R}$）的部分对应值如下表：

\begin{center}
\begingroup\setlength{\tabcolsep}{4pt}
\begin{tabular}{ccccccccc}
\toprule
$x$ & $-3$ & $-2$ & $-1$ & $0$ & $1$ & $2$ & $3$ & $4$ \\
\midrule
$y$ & $6$ & $0$ & $-4$ & $-6$ & $-6$ & $-4$ & $0$ & $6$ \\
\bottomrule
\end{tabular}
\endgroup
\end{center}

则不等式 $ax^2 + bx + c > 0$ 的解集是 \underline{\hspace{3cm}}．

\Needspace{4\baselineskip}
\textbf{解析}

由题给表格可知：
当 $x = -2$ 和 $x = 3$ 时，函数值均为 $y = 0$，故方程 $ax^2 + bx + c = 0$ 的两个根为 $x_1 = -2, x_2 = 3$．
且由表可知，当 $x < -2$ 或 $x > 3$ 时，函数值 $y > 0$；
当 $x = 0$ 时，$y = c = -6$．由 $a(0 - (-2))(0 - 3) = -6a = -6$ 得二次项系数 $a = 1 > 0$．
因此不等式 $ax^2 + bx + c > 0$ 即 $(x + 2)(x - 3) > 0$．
其解集为 $\{x \mid x < -2 \text{ \text{或} } x > 3\}$（或记为 $(-\infty, -2) \cup (3, +\infty)$）．

\gkanswer{答案：$\{x \mid x < -2 \text{ \text{或} } x > 3\}$（或 $(-\infty, -2) \cup (3, +\infty)$）}

\end{answer}
\begin{answer}
\gksolhead{14}{}

\textbf{题目}：

以点 $(1, 2)$ 为圆心，与直线 $4x + 3y - 35 = 0$ 相切的圆的方程是 \underline{\hspace{3cm}}．

\Needspace{4\baselineskip}
\textbf{解析}

圆心为 $(1, 2)$，因为该圆与直线 $4x + 3y - 35 = 0$ 相切，
\Needspace{5\baselineskip}
所以圆的半径 $r$ 等于圆心到直线的点到直线距离：
\nopagebreak[4]
\[
r = \frac{|4 \times 1 + 3 \times 2 - 35|}{\sqrt{4^2 + 3^2}} = \frac{|4 + 6 - 35|}{5} = \frac{|-25|}{5} = 5.
\]
\Needspace{5\baselineskip}
因此所求圆的标准方程为
\nopagebreak[4]
\[
(x - 1)^2 + (y - 2)^2 = 25.
\]

\gkanswer{答案：$(x - 1)^2 + (y - 2)^2 = 25$}

\end{answer}
\begin{answer}
\gksolhead{15}{}

\textbf{题目}：

设数列 $\{a_n\}$ 的前 $n$ 项和为 $S_n$，$S_n = \dfrac{a_1(3^n - 1)}{2}$（对于所有 $n \geqslant 1$），且 $a_4 = 54$，则 $a_1$ 的数值是 \underline{\hspace{3cm}}．

\Needspace{4\baselineskip}
\textbf{解析}

已知数列 $\{a_n\}$ 的前 $n$ 项和为 $S_n = \dfrac{a_1(3^n - 1)}{2}$．
\Needspace{5\baselineskip}
当 $n \geqslant 2$ 时，
\nopagebreak[4]
\[
a_n = S_n - S_{n-1} = \frac{a_1(3^n - 1)}{2} - \frac{a_1(3^{n-1} - 1)}{2} = \frac{a_1}{2} \left(3^n - 3^{n-1}\right) = \frac{a_1}{2} \cdot 3^{n-1}(3 - 1) = a_1 \cdot 3^{n-1}.
\]
当 $n = 1$ 时，$S_1 = \dfrac{a_1(3^1 - 1)}{2} = a_1$，通项公式对 $n = 1$ 也成立．
因此 $\{a_n\}$ 是公比为 $q = 3$ 的等比数列．
\Needspace{5\baselineskip}
由 $a_4 = 54$，得
\nopagebreak[4]
\[
a_1 \cdot 3^{4-1} = 27 a_1 = 54 \implies a_1 = 2.
\]

\gkanswer{答案：$2$}

\end{answer}
\begin{answer}
\gksolhead{16}{}

\textbf{题目}：

平面向量 $\vec{a}, \vec{b}$ 中，已知 $\vec{a} = (4, -3)$，$|\vec{b}| = 1$，且 $\vec{a} \cdot \vec{b} = 5$，则向量 $\vec{b} =$ \underline{\hspace{3cm}}．

\Needspace{4\baselineskip}
\textbf{解析}

\Needspace{5\baselineskip}
已知 $\vec{a} = (4, -3)$，则其模长为
\nopagebreak[4]
\[
|\vec{a}| = \sqrt{4^2 + (-3)^2} = 5.
\]
已知 $|\vec{b}| = 1$，设向量 $\vec{a}$ 与 $\vec{b}$ 的夹角为 $\theta$．
\Needspace{5\baselineskip}
由数量积定义：
\nopagebreak[4]
\[
\vec{a} \cdot \vec{b} = |\vec{a}| |\vec{b}| \cos\theta = 5 \times 1 \times \cos\theta = 5\cos\theta.
\]
已知 $\vec{a} \cdot \vec{b} = 5$，故 $\cos\theta = 1$，这表明向量 $\vec{b}$ 与 $\vec{a}$ 同向共线．
\Needspace{5\baselineskip}
因此
\nopagebreak[4]
\[
\vec{b} = \frac{\vec{a}}{|\vec{a}|} = \frac{(4, -3)}{5} = \left(\frac{4}{5}, -\frac{3}{5}\right).
\]

\gkanswer{答案：$\left(\dfrac{4}{5}, -\dfrac{3}{5}\right)$}

\end{answer}
\begin{answer}
\gksolhead{17}{}

\textbf{题目}（本小题满分 12 分）：

已知 $0 < \alpha < \dfrac{\pi}{2}$，$\tan\dfrac{\alpha}{2} + \cot\dfrac{\alpha}{2} = \dfrac{5}{2}$，求 $\sin\left(\alpha - \dfrac{\pi}{3}\right)$ 的值．

\Needspace{4\baselineskip}
\textbf{解析}

\Needspace{5\baselineskip}
利用切化弦化简已知条件：
\nopagebreak[4]
\[
\tan\frac{\alpha}{2} + \cot\frac{\alpha}{2} = \frac{\sin\frac{\alpha}{2}}{\cos\frac{\alpha}{2}} + \frac{\cos\frac{\alpha}{2}}{\sin\frac{\alpha}{2}} = \frac{\sin^2\frac{\alpha}{2} + \cos^2\frac{\alpha}{2}}{\sin\frac{\alpha}{2}\cos\frac{\alpha}{2}} = \frac{1}{\frac{1}{2}\sin\alpha} = \frac{2}{\sin\alpha}.
\]
因为 $\tan\dfrac{\alpha}{2} + \cot\dfrac{\alpha}{2} = \dfrac{5}{2}$，
\Needspace{5\baselineskip}
所以
\nopagebreak[4]
\[
\frac{2}{\sin\alpha} = \frac{5}{2} \implies \sin\alpha = \frac{4}{5}.
\]
\Needspace{5\baselineskip}
已知 $0 < \alpha < \dfrac{\pi}{2}$，所以
\nopagebreak[4]
\[
\cos\alpha = \sqrt{1 - \sin^2\alpha} = \sqrt{1 - \left(\frac{4}{5}\right)^2} = \frac{3}{5}.
\]
\Needspace{5\baselineskip}
利用两角差的正弦公式展开：
\nopagebreak[4]
\[
\sin\left(\alpha - \frac{\pi}{3}\right) = \sin\alpha \cos\frac{\pi}{3} - \cos\alpha \sin\frac{\pi}{3} = \frac{4}{5} \times \frac{1}{2} - \frac{3}{5} \times \frac{\sqrt{3}}{2} = \frac{4 - 3\sqrt{3}}{10}.
\]

\gkanswer{答案：$\sin\left(\alpha - \dfrac{\pi}{3}\right) = \dfrac{4 - 3\sqrt{3}}{10}$}

\end{answer}
\begin{answer}
\gksolhead{18}{}

\textbf{题目}（本小题满分 12 分）：

在棱长为 4 的正方体 $ABCD - A_1B_1C_1D_1$ 中，$O$ 是正方形 $A_1B_1C_1D_1$ 的中心，点 $P$ 在棱 $CC_1$ 上，且 $CC_1 = 4CP$．

（1）求直线 $AP$ 与平面 $BCC_1B_1$ 所成的角的大小（结果用反三角函数值表示）；

（2）设 $O$ 点在平面 $D_1AP$ 上的射影是 $H$，求证：$D_1H \perp AP$；

（3）求点 $P$ 到平面 $ABD_1$ 的距离．
\nopagebreak[4]
\begin{center}\begin{tikzpicture}[>=stealth, scale=1.0]
  % 顶点坐标
  \coordinate (D) at (0, 0);
  \coordinate (C) at (3.2, 0);
  \coordinate (A) at (-1.5, -1.2);
  \coordinate (B) at (1.7, -1.2);

  \coordinate (D1) at (0, 3.0);
  \coordinate (C1) at (3.2, 3.0);
  \coordinate (A1) at (-1.5, 1.8);
  \coordinate (B1) at (1.7, 1.8);

  % 点 P 在 CC1 上，CP = 1/4 CC1
  \coordinate (P) at ($(C)!0.25!(C1)$);

  % 上底面中心 O
  \coordinate (O) at ($(A1)!0.5!(C1)$);

  % 射影点 H (在平面 D1AP 上)
  \coordinate (H) at (0.9, 1.4);

  % 看不见的棱 (虚线)
  \draw[dashed, thick] (D) -- (A);
  \draw[dashed, thick] (D) -- (C);
  \draw[dashed, thick] (D) -- (D1);

  % 外轮廓与可见棱 (实线)
  \draw[thick] (A) -- (B) -- (C);
  \draw[thick] (A1) -- (B1) -- (C1) -- (D1) -- cycle;
  \draw[thick] (A) -- (A1);
  \draw[thick] (B) -- (B1);
  \draw[thick] (C) -- (C1);

  % 连线 AP, D1A, D1P
  \draw[thick, blue!80!black] (D1) -- (A);
  \draw[thick, blue!80!black] (A) -- (P);
  \draw[thick, blue!80!black] (D1) -- (P);

  % 辅助线 D1O, OH, D1H
  \draw[thick] (D1) -- (O);
  \draw[dashed] (O) -- (H);
  \draw[thick, red!70!black] (D1) -- (H);

  % 标记点
  \fill (O) circle (1.5pt) node[above right] {$O$};
  \fill (P) circle (1.5pt) node[right] {$P$};
  \fill (H) circle (1.5pt) node[below right] {$H$};

  % 顶点标签
  \node[left] at (A) {$A$};
  \node[below right] at (B) {$B$};
  \node[right] at (C) {$C$};
  \node[above left] at (D) {$D$};
  \node[left] at (A1) {$A_1$};
  \node[below right] at (B1) {$B_1$};
  \node[above right] at (C1) {$C_1$};
  \node[above left] at (D1) {$D_1$};
\end{tikzpicture}
\end{center}
\nopagebreak[4]
\Needspace{4\baselineskip}
\textbf{解析}
\nopagebreak[4]
\begin{center}\begin{tikzpicture}[>=stealth, scale=1.0]
  % 顶点坐标
  \coordinate (D) at (0, 0);
  \coordinate (C) at (3.2, 0);
  \coordinate (A) at (-1.5, -1.2);
  \coordinate (B) at (1.7, -1.2);

  \coordinate (D1) at (0, 3.0);
  \coordinate (C1) at (3.2, 3.0);
  \coordinate (A1) at (-1.5, 1.8);
  \coordinate (B1) at (1.7, 1.8);

  % 点 P 在 CC1 上，CP = 1/4 CC1
  \coordinate (P) at ($(C)!0.25!(C1)$);

  % 上底面中心 O
  \coordinate (O) at ($(A1)!0.5!(C1)$);

  % 射影点 H (在平面 D1AP 上)
  \coordinate (H) at (0.9, 1.4);

  % 看不见的棱 (虚线)
  \draw[dashed, thick] (D) -- (A);
  \draw[dashed, thick] (D) -- (C);
  \draw[dashed, thick] (D) -- (D1);

  % 外轮廓与可见棱 (实线)
  \draw[thick] (A) -- (B) -- (C);
  \draw[thick] (A1) -- (B1) -- (C1) -- (D1) -- cycle;
  \draw[thick] (A) -- (A1);
  \draw[thick] (B) -- (B1);
  \draw[thick] (C) -- (C1);

  % 连线 AP, D1A, D1P
  \draw[thick, blue!80!black] (D1) -- (A);
  \draw[thick, blue!80!black] (A) -- (P);
  \draw[thick, blue!80!black] (D1) -- (P);

  % 辅助线 D1O, OH, D1H
  \draw[thick] (D1) -- (O);
  \draw[dashed] (O) -- (H);
  \draw[thick, red!70!black] (D1) -- (H);

  % 标记点
  \fill (O) circle (1.5pt) node[above right] {$O$};
  \fill (P) circle (1.5pt) node[right] {$P$};
  \fill (H) circle (1.5pt) node[below right] {$H$};

  % 顶点标签
  \node[left] at (A) {$A$};
  \node[below right] at (B) {$B$};
  \node[right] at (C) {$C$};
  \node[above left] at (D) {$D$};
  \node[left] at (A1) {$A_1$};
  \node[below right] at (B1) {$B_1$};
  \node[above right] at (C1) {$C_1$};
  \node[above left] at (D1) {$D_1$};
\end{tikzpicture}
\end{center}
\nopagebreak[4]
（1）在正方体 $ABCD - A_1B_1C_1D_1$ 中，
因为 $AB \perp$ 平面 $BCC_1B_1$，
所以点 $A$ 在平面 $BCC_1B_1$ 上的射影为点 $B$．
连接 $BP$，则 $BP$ 为直线 $AP$ 在平面 $BCC_1B_1$ 上的射影，
因此 $\angle APB$ 就是直线 $AP$ 与平面 $BCC_1B_1$ 所成的角．
在 Rt$\triangle ABP$ 中，$AB \perp BP$，$AB = 4$．
在正方形 $BCC_1B_1$ 中，$BC = 4$，$CP = \dfrac{1}{4}CC_1 = 1$，
\Needspace{5\baselineskip}
所以
\nopagebreak[4]
\[
BP = \sqrt{BC^2 + CP^2} = \sqrt{4^2 + 1^2} = \sqrt{17}.
\]
\Needspace{5\baselineskip}
在 Rt$\triangle ABP$ 中，
\nopagebreak[4]
\[
\tan\angle APB = \frac{AB}{BP} = \frac{4}{\sqrt{17}} = \frac{4\sqrt{17}}{17}.
\]
所以直线 $AP$ 与平面 $BCC_1B_1$ 所成的角为 $\arctan\dfrac{4\sqrt{17}}{17}$（或 $\arctan\dfrac{4}{\sqrt{17}}$）．

（2）以点 $D$ 为原点，$DA, DC, DD_1$ 所在直线分别为 $x, y, z$ 轴建立空间直角坐标系．
正方体棱长为 4，各点坐标为：
$D(0, 0, 0)$，$A(4, 0, 0)$，$B(4, 4, 0)$，$C(0, 4, 0)$，
$D_1(0, 0, 4)$，$A_1(4, 0, 4)$，$B_1(4, 4, 4)$，$C_1(0, 4, 4)$．
$O$ 是正方形 $A_1B_1C_1D_1$ 的中心，故 $O(2, 2, 4)$．
因为 $CC_1 = 4CP$，$CP = 1$，所以点 $P(0, 4, 1)$．
\Needspace{5\baselineskip}
计算向量坐标：
\nopagebreak[4]
\[
\vec{AP} = P - A = (-4, 4, 1),
\]
\nopagebreak[4]
\[
\vec{D_1O} = O - D_1 = (2, 2, 0).
\]
\Needspace{5\baselineskip}
计算数量积：
\nopagebreak[4]
\[
\vec{D_1O} \cdot \vec{AP} = 2 \times (-4) + 2 \times 4 + 0 \times 1 = -8 + 8 + 0 = 0.
\]
所以 $D_1O \perp AP$．
又已知点 $H$ 是 $O$ 在平面 $D_1AP$ 上的射影，故 $OH \perp$ 平面 $D_1AP$．
因为 $AP \subset$ 平面 $D_1AP$，所以 $OH \perp AP$．
又 $D_1O \cap OH = O$，故 $AP \perp$ 平面 $D_1OH$．
因为 $D_1H \subset$ 平面 $D_1OH$，所以 $D_1H \perp AP$．

（3）连接 $AD_1, BD_1, AB$．在空间直角坐标系中，
$A(4, 0, 0)$，$B(4, 4, 0)$，$D_1(0, 0, 4)$．
向量 $\vec{AB} = (0, 4, 0)$，$\vec{AD_1} = (-4, 0, 4)$．
\Needspace{5\baselineskip}
设平面 $ABD_1$ 的法向量为 $\vec{n} = (x, y, z)$，则
\nopagebreak[4]
\[
\begin{cases}
\vec{n} \cdot \vec{AB} = 4y = 0, \\
\vec{n} \cdot \vec{AD_1} = -4x + 4z = 0.
\end{cases}
\]
取 $x = 1$，则 $y = 0, z = 1$，得 $\vec{n} = (1, 0, 1)$．
由（2）知 $\vec{AP} = (-4, 4, 1)$，
\Needspace{5\baselineskip}
所以点 $P$ 到平面 $ABD_1$ 的距离为
\nopagebreak[4]
\[
d = \frac{|\vec{AP} \cdot \vec{n}|}{|\vec{n}|} = \frac{|-4 \times 1 + 4 \times 0 + 1 \times 1|}{\sqrt{1^2 + 0^2 + 1^2}} = \frac{|-3|}{\sqrt{2}} = \frac{3\sqrt{2}}{2}.
\]

\gkanswer{答案：（1）直线 $AP$ 与平面 $BCC_1B_1$ 所成的角的大小为 $\arctan\dfrac{4\sqrt{17}}{17}$（或 $\arctan\dfrac{4}{\sqrt{17}}$）；（2）证明见解析；（3）点 $P$ 到平面 $ABD_1$ 的距离为 $\dfrac{3\sqrt{2}}{2}$}

\end{answer}
\begin{answer}
\gksolhead{19}{}

\textbf{题目}（本小题满分 12 分）：

制定投资计划时，不仅要考虑可能获得的盈利，而且要考虑可能出现的亏损．某投资人打算投资甲、乙两个项目．根据预测，甲、乙项目可能的最大盈利率分别为 $100\%$ 和 $50\%$，可能的最大亏损分别为 $30\%$ 和 $10\%$．投资人计划投资金额不超过 10 万元，要求确保可能的资金亏损不超过 1.8 万元．问投资人对甲、乙两个项目各投资多少万元，才能使可能的盈利最大？

\Needspace{4\baselineskip}
\textbf{解析}

设投资人对甲项目投资 $x$ 万元，对乙项目投资 $y$ 万元，获得的可能盈利为 $z$ 万元．
由题意，投资总金额不超过 10 万元，且确保可能资金亏损不超过 1.8 万元，
\Needspace{5\baselineskip}
所以 $x, y$ 满足的线性约束条件为
\nopagebreak[4]
\[
\begin{cases}
x + y \leqslant 10, \\
30\% x + 10\% y \leqslant 1.8, \\
x \geqslant 0, \\
y \geqslant 0,
\end{cases}
\quad \text{\text{即}} \quad
\begin{cases}
x + y \leqslant 10, \\
3x + y \leqslant 18, \\
x \geqslant 0, \\
y \geqslant 0.
\end{cases}
\]
\Needspace{5\baselineskip}
目标函数为可能的最大盈利：
\nopagebreak[4]
\[
z = 100\% x + 50\% y = x + 0.5y.
\]
可行域为平面直角坐标系中由四条直线围成的四边形区域，其顶点坐标分别为：
\begin{itemize}
\item $O(0, 0)$；
\item $A(6, 0)$（直线 $3x + y = 18$ 与 $x$ 轴交点）；
\item $B(4, 6)$（联立 $x + y = 10$ 与 $3x + y = 18$ 解得）；
\item $C(0, 10)$（直线 $x + y = 10$ 与 $y$ 轴交点）．
\end{itemize}

分别计算各顶点处的盈利值 $z$：
\begin{itemize}
\item 在点 $O(0, 0)$ 处，$z = 0$；
\item 在点 $A(6, 0)$ 处，$z = 6 + 0.5 \times 0 = 6$ 万元；
\item 在点 $B(4, 6)$ 处，$z = 4 + 0.5 \times 6 = 7$ 万元；
\item 在点 $C(0, 10)$ 处，$z = 0 + 0.5 \times 10 = 5$ 万元．
\end{itemize}

比较可知，当 $x = 4, y = 6$ 时，$z$ 取得最大值 7 万元．
答：投资人对甲项目投资 4 万元，对乙项目投资 6 万元，才能使可能的盈利最大，最大可能盈利为 7 万元．

\gkanswer{答案：投资人对甲项目投资 4 万元，对乙项目投资 6 万元，才能使可能的盈利最大，最大盈利为 7 万元}

\end{answer}
\begin{answer}
\gksolhead{20}{}

\textbf{题目}（本小题满分 12 分）：

设无穷等差数列 $\{a_n\}$ 的前 $n$ 项和为 $S_n$．

（1）若首项 $a_1 = \dfrac{3}{2}$，公差 $d = 1$，求满足 $S_{k^2} = (S_k)^2$ 的正整数 $k$；

（2）求所有的无穷等差数列 $\{a_n\}$，使得对于一切正整数 $k$ 都有 $S_{k^2} = (S_k)^2$ 成立．

\Needspace{4\baselineskip}
\textbf{解析}

\Needspace{5\baselineskip}
（1）当 $a_1 = \dfrac{3}{2}, d = 1$ 时，等差数列前 $n$ 项和为
\nopagebreak[4]
\[
S_n = n a_1 + \frac{n(n - 1)}{2} d = \frac{3}{2}n + \frac{n(n - 1)}{2} = \frac{n^2 + 2n}{2} = \frac{n(n + 2)}{2}.
\]
\Needspace{5\baselineskip}
因此
\nopagebreak[4]
\[
S_{k^2} = \frac{k^2(k^2 + 2)}{2},
\]
\nopagebreak[4]
\[
(S_k)^2 = \left[\frac{k(k + 2)}{2}\right]^2 = \frac{k^2(k + 2)^2}{4}.
\]
\Needspace{5\baselineskip}
由 $S_{k^2} = (S_k)^2$，两边同除以 $\dfrac{k^2}{4}$（因为 $k$ 是正整数，$k \neq 0$）：
\nopagebreak[4]
\[
2(k^2 + 2) = (k + 2)^2 \implies 2k^2 + 4 = k^2 + 4k + 4 \implies k^2 - 4k = 0.
\]
因为 $k$ 是正整数，解得 $k = 4$．

\Needspace{5\baselineskip}
（2）设无穷等差数列的前 $n$ 项和为
\nopagebreak[4]
\[
S_n = \frac{d}{2}n^2 + \left(a_1 - \frac{d}{2}\right)n.
\]
记 $A = \dfrac{d}{2}, B = a_1 - \dfrac{d}{2}$，则 $S_n = A n^2 + B n$．
\Needspace{5\baselineskip}
对于一切正整数 $k$ 都有 $S_{k^2} = (S_k)^2$ 成立，即
\nopagebreak[4]
\[
A(k^2)^2 + B(k^2) = (A k^2 + B k)^2,
\]
\nopagebreak[4]
\[
A k^4 + B k^2 = A^2 k^4 + 2AB k^3 + B^2 k^2.
\]
\Needspace{5\baselineskip}
该等式对于一切正整数 $k$ 恒成立，因此对应项的系数必须恒等：
\nopagebreak[4]
\[
\begin{cases}
A^2 = A, \\
2AB = 0, \\
B^2 = B.
\end{cases}
\]
由 $A^2 = A$ 得 $A = 0$ 或 $A = 1$；
由 $B^2 = B$ 得 $B = 0$ 或 $B = 1$；
结合 $2AB = 0$ 知 $A$ 与 $B$ 不能同时为 1．
因此只有以下三种解的情况：
1. $A = 0, B = 0$：
   $\dfrac{d}{2} = 0 \implies d = 0$，$a_1 - 0 = 0 \implies a_1 = 0$．
   此时数列为 $a_n = 0$（对一切 $n \in \mathbf{N}^*$）．
2. $A = 1, B = 0$：
   $\dfrac{d}{2} = 1 \implies d = 2$，$a_1 - 1 = 0 \implies a_1 = 1$．
   此时数列通项为 $a_n = 1 + (n - 1) \times 2 = 2n - 1$．
3. $A = 0, B = 1$：
   $\dfrac{d}{2} = 0 \implies d = 0$，$a_1 - 0 = 1 \implies a_1 = 1$．
   此时数列为 $a_n = 1$（对一切 $n \in \mathbf{N}^*$）．

综上所述，满足条件的一切无穷等差数列 $\{a_n\}$ 共有三个：
$a_n = 0$；$a_n = 1$；$a_n = 2n - 1$．

\gkanswer{答案：（1）$k = 4$；（2）满足条件的无穷等差数列 $\{a_n\}$ 共有三个：$a_n = 0$；$a_n = 1$；$a_n = 2n - 1$}

\end{answer}
\begin{answer}
\gksolhead{21}{}

\textbf{题目}（本小题满分 12 分）：

已知椭圆的中心在原点，离心率为 $\dfrac{1}{2}$，一个焦点是 $F(-m, 0)$（$m$ 是大于 0 的常数）．

（1）求椭圆的方程；

（2）设 $Q$ 是椭圆上的一点，且过点 $F, Q$ 的直线 $l$ 与 $y$ 轴交于点 $M$．若 $\vec{MQ} = 2\vec{QF}$，求直线 $l$ 的斜率．

\Needspace{4\baselineskip}
\textbf{解析}

（1）设椭圆方程为 $\dfrac{x^2}{a^2} + \dfrac{y^2}{b^2} = 1$（$a > b > 0$）．
由焦点 $F(-m, 0)$ 可知 $c = m$．
离心率 $e = \dfrac{c}{a} = \dfrac{1}{2}$，所以 $a = 2c = 2m$．
\Needspace{5\baselineskip}
由椭圆中参数关系 $b^2 = a^2 - c^2$ 得
\nopagebreak[4]
\[
b^2 = (2m)^2 - m^2 = 3m^2.
\]
\Needspace{5\baselineskip}
因此所求椭圆的标准方程为
\nopagebreak[4]
\[
\frac{x^2}{4m^2} + \frac{y^2}{3m^2} = 1.
\]

（2）直线 $l$ 过焦点 $F(-m, 0)$ 且与 $y$ 轴交于点 $M$，故直线 $l$ 不垂直于 $x$ 轴．
\Needspace{5\baselineskip}
设直线 $l$ 的斜率为 $k$，则直线方程为
\nopagebreak[4]
\[
y = k(x + m).
\]
令 $x = 0$，得交点 $M(0, km)$．
设点 $Q$ 的坐标为 $(x_0, y_0)$．
\Needspace{5\baselineskip}
则向量
\nopagebreak[4]
\[
\vec{MQ} = (x_0 - 0, y_0 - km) = (x_0, y_0 - km),
\]
\nopagebreak[4]
\[
\vec{QF} = (-m - x_0, 0 - y_0) = (-m - x_0, -y_0).
\]
\Needspace{5\baselineskip}
由已知条件 $\vec{MQ} = 2\vec{QF}$，得
\nopagebreak[4]
\[
\begin{cases}
x_0 = 2(-m - x_0), \\
y_0 - km = 2(-y_0).
\end{cases}
\]
由第一式解得 $x_0 = -2m - 2x_0 \implies 3x_0 = -2m \implies x_0 = -\dfrac{2}{3}m$；
由第二式解得 $3y_0 = km \implies y_0 = \dfrac{1}{3}km$．
\Needspace{5\baselineskip}
因为点 $Q(x_0, y_0)$ 在椭圆上，将其代入椭圆方程：
\nopagebreak[4]
\[
\frac{\left(-\frac{2}{3}m\right)^2}{4m^2} + \frac{\left(\frac{1}{3}km\right)^2}{3m^2} = 1,
\]
\nopagebreak[4]
\[
\frac{\frac{4}{9}m^2}{4m^2} + \frac{\frac{1}{9}k^2 m^2}{3m^2} = 1 \implies \frac{1}{9} + \frac{k^2}{27} = 1.
\]
\Needspace{5\baselineskip}
两边同乘以 27：
\nopagebreak[4]
\[
3 + k^2 = 27 \implies k^2 = 24 \implies k = \pm 2\sqrt{6}.
\]
因此直线 $l$ 的斜率为 $\pm 2\sqrt{6}$．

\gkanswer{答案：（1）椭圆的方程为 $\dfrac{x^2}{4m^2} + \dfrac{y^2}{3m^2} = 1$；（2）直线 $l$ 的斜率为 $\pm 2\sqrt{6}$}

\end{answer}
\begin{answer}
\gksolhead{22}{}

\textbf{题目}（本小题满分 14 分）：

\Needspace{5\baselineskip}
已知函数 $f(x)$（$x \in \mathbf{R}$）满足下列条件：对任意的实数 $x_1, x_2$ 都有
\nopagebreak[4]
\[
\lambda(x_1 - x_2)^2 \leqslant (x_1 - x_2)[f(x_1) - f(x_2)] \quad \text{\text{和}} \quad |f(x_1) - f(x_2)| \leqslant |x_1 - x_2|,
\]
其中 $\lambda$ 是大于 0 的常数．设实数 $a_0, a, b$ 满足 $f(a_0) = 0$ 和 $b = a - \lambda f(a)$．

（1）证明：$\lambda \leqslant 1$，并且不存在 $b_0 \neq a_0$，使得 $f(b_0) = 0$；

（2）证明：$(b - a_0)^2 \leqslant (1 - \lambda^2)(a - a_0)^2$；

（3）证明：$[f(b)]^2 \leqslant (1 - \lambda^2)[f(a)]^2$．

\Needspace{4\baselineskip}
\textbf{解析}

（1）证明：
\Needspace{5\baselineskip}
对任意两个不相等的实数 $x_1 \neq x_2$，由已知条件：
\nopagebreak[4]
\[
\lambda(x_1 - x_2)^2 \leqslant (x_1 - x_2)[f(x_1) - f(x_2)] \leqslant |x_1 - x_2||f(x_1) - f(x_2)|.
\]
\Needspace{5\baselineskip}
结合第二个条件 $|f(x_1) - f(x_2)| \leqslant |x_1 - x_2|$，可得
\nopagebreak[4]
\[
\lambda(x_1 - x_2)^2 \leqslant |x_1 - x_2| \cdot |x_1 - x_2| = (x_1 - x_2)^2.
\]
因为 $x_1 \neq x_2$，所以 $(x_1 - x_2)^2 > 0$．
\Needspace{5\baselineskip}
两边同除以 $(x_1 - x_2)^2$，即得
\nopagebreak[4]
\[
\lambda \leqslant 1.
\]
假设存在实数 $b_0 \neq a_0$ 使得 $f(b_0) = 0$．
\Needspace{5\baselineskip}
在第一个已知条件中取 $x_1 = a_0, x_2 = b_0$，则
\nopagebreak[4]
\[
\lambda(a_0 - b_0)^2 \leqslant (a_0 - b_0)[f(a_0) - f(b_0)] = (a_0 - b_0)(0 - 0) = 0.
\]
因为 $\lambda > 0$，所以必有 $(a_0 - b_0)^2 \leqslant 0$．
从而只能有 $a_0 = b_0$，这与假设 $b_0 \neq a_0$ 矛盾．
因此不存在 $b_0 \neq a_0$ 使得 $f(b_0) = 0$．

（2）证明：
\Needspace{5\baselineskip}
由 $b = a - \lambda f(a)$，得
\nopagebreak[4]
\[
b - a_0 = (a - a_0) - \lambda f(a).
\]
\Needspace{5\baselineskip}
两边平方展开得
\nopagebreak[4]
\[
(b - a_0)^2 = (a - a_0)^2 - 2\lambda(a - a_0)f(a) + \lambda^2 [f(a)]^2.
\]
在已知条件 $\lambda(x_1 - x_2)^2 \leqslant (x_1 - x_2)[f(x_1) - f(x_2)]$ 中，取 $x_1 = a, x_2 = a_0$．
\Needspace{5\baselineskip}
因为 $f(a_0) = 0$，所以
\nopagebreak[4]
\[
\lambda(a - a_0)^2 \leqslant (a - a_0)[f(a) - f(a_0)] = (a - a_0)f(a).
\]
\Needspace{5\baselineskip}
两边同乘以 $-2\lambda$（因为 $\lambda > 0$，故 $-2\lambda < 0$），不等号反向：
\nopagebreak[4]
\[
-2\lambda(a - a_0)f(a) \leqslant -2\lambda^2(a - a_0)^2. \tag{*}
\]
\Needspace{5\baselineskip}
又由第二个已知条件，取 $x_1 = a, x_2 = a_0$：
\nopagebreak[4]
\[
|f(a)| = |f(a) - f(a_0)| \leqslant |a - a_0|.
\]
两边平方得 $[f(a)]^2 \leqslant (a - a_0)^2$．
\Needspace{5\baselineskip}
两边同乘以 $\lambda^2 > 0$，得
\nopagebreak[4]
\[
\lambda^2 [f(a)]^2 \leqslant \lambda^2 (a - a_0)^2. \tag{**}
\]
\Needspace{5\baselineskip}
将不等式 (*) 和 (**) 代入 $(b - a_0)^2$ 的展开式，得
\nopagebreak[4]
\[
(b - a_0)^2 \leqslant (a - a_0)^2 - 2\lambda^2(a - a_0)^2 + \lambda^2(a - a_0)^2 = (1 - \lambda^2)(a - a_0)^2.
\]
命题得证．

（3）证明：
\Needspace{5\baselineskip}
将 $f(b)$ 变形为
\nopagebreak[4]
\[
f(b) = f(a) + [f(b) - f(a)].
\]
\Needspace{5\baselineskip}
两边平方展开得
\nopagebreak[4]
\[
[f(b)]^2 = [f(a)]^2 + 2f(a)[f(b) - f(a)] + [f(b) - f(a)]^2.
\]
\Needspace{5\baselineskip}
因为 $b - a = -\lambda f(a)$，在已知条件 $\lambda(x_1 - x_2)^2 \leqslant (x_1 - x_2)[f(x_1) - f(x_2)]$ 中，取 $x_1 = b, x_2 = a$，得
\nopagebreak[4]
\[
\lambda(b - a)^2 \leqslant (b - a)[f(b) - f(a)].
\]
\Needspace{5\baselineskip}
将 $b - a = -\lambda f(a)$ 代入上式：
\nopagebreak[4]
\[
\lambda[-\lambda f(a)]^2 \leqslant -\lambda f(a)[f(b) - f(a)] \implies \lambda^3 [f(a)]^2 \leqslant -\lambda f(a)[f(b) - f(a)].
\]
\Needspace{5\baselineskip}
两边同除以 $\lambda > 0$，得
\nopagebreak[4]
\[
\lambda^2 [f(a)]^2 \leqslant -f(a)[f(b) - f(a)] \implies f(a)[f(b) - f(a)] \leqslant -\lambda^2 [f(a)]^2.
\]
\Needspace{5\baselineskip}
两边同乘以 2，得
\nopagebreak[4]
\[
2f(a)[f(b) - f(a)] \leqslant -2\lambda^2 [f(a)]^2. \tag{1}
\]
\Needspace{5\baselineskip}
又由第二个已知条件，取 $x_1 = b, x_2 = a$：
\nopagebreak[4]
\[
|f(b) - f(a)| \leqslant |b - a| = |-\lambda f(a)| = \lambda |f(a)|.
\]
\Needspace{5\baselineskip}
两边平方得
\nopagebreak[4]
\[
[f(b) - f(a)]^2 \leqslant \lambda^2 [f(a)]^2. \tag{2}
\]
\Needspace{5\baselineskip}
将不等式 (1) 与 (2) 代入 $[f(b)]^2$ 的展开式中：
\nopagebreak[4]
\[
[f(b)]^2 \leqslant [f(a)]^2 - 2\lambda^2 [f(a)]^2 + \lambda^2 [f(a)]^2 = (1 - \lambda^2)[f(a)]^2.
\]
命题得证．

\gkanswer{答案：（1）证明见解析；（2）证明见解析；（3）证明见解析}

\end{answer}